% theory/varcurv/varcurv.tex -- variable curvature: what survives of the power-sum structure.
% LaTeX fragment (amsthm environments theorem/lemma/proposition/corollary/remark assumed).
% Scripts: twisted_mp.py, top_coefficient.py, quadratic_form.py, obstruction_check.py (all asserting).

\section*{Cone points in variable curvature}

\paragraph{Setting.}
$(\Orb,g)$ is a closed orientable Riemannian $2$-orbifold whose singular points are cone points
$p_1,\dots,p_n$ of orders $m_i\ge2$, on an underlying surface $|\Orb|$ of genus $\gamma$, and
$\chi^{\rm orb}=2-2\gamma-\sum_i(1-1/m_i)$. $K$ is the Gauss curvature and $\Delta_g=-\operatorname{div}\operatorname{grad}$.
By Donnelly and Dryden--Gordon--Greenwald--Webb \cite[4.1, Thm~4.8, Prop.~4.11]{dggw2008},
\begin{equation}\label{vc:exp}
\operatorname{tr}e^{-t\Delta_g}\sim\sum_{l\ge-1}a_l\,t^l,\qquad
a_l=S_l+\sum_{i=1}^na_l(p_i),\qquad S_l=\frac1{4\pi}\int_\Orb u_{l+1}(x,x)\,dA,
\end{equation}
with the Minakshisundaram--Pleijel coefficients $u_0=1$, $u_1=K/3$, $u_2=K^2/15-\Delta_gK/15$
\cite[(7)]{schueth2019}, and $a_l(p)=\frac1m\sum_{j=1}^{m-1}b_l(\Phi^j)$ \cite[(17)]{schueth2019},
$\Phi$ a rotation of order $m$ in an orbifold chart and $b_l$ Donnelly's coefficient of the isolated
fixed point. Paper~A's $c_j$ is $a_{j-2}$. Write $C_\phi^2=2-2\cos\phi=4\sin^2(\phi/2)$ and
\[
\Pi_i(m)=\frac1m\sum_{j=1}^{m-1}\bigl(4\sin^2(\pi j/m)\bigr)^{-i}.
\]
$S_{-1}=\Area/4\pi$ and $S_0=\chi^{\rm orb}/6$ (Gauss--Bonnet); for $l\ge1$, $S_l$ is the integral of a
universal polynomial in $K$ and its covariant derivatives, each monomial of degree $\ge1$, so $S_l$
vanishes on flat regions. Weights: $\nabla^jK$ has weight $j+2$; $a_l(p)$ and $u_{l+1}$ have weight $2l$
and $2l+2$; a monomial of weight $2l$ involves $\nabla^jK$ only for $j\le2l-2$.

\subsection*{1. What is known in closed form}

\begin{center}\small
\begin{tabular}{llll}
order & cone point & jets entering & source\\\hline
$t^0$ & orbifold, order $m$, any metric & none: $\frac1{12}(m-\frac1m)$ & \cite[Prop.~5.5]{dggw2008}\\
$t^1$ & orbifold, any metric & $K(p)$: $[\frac1{360}(m^3-\frac1m)+\frac1{36}(m-\frac1m)]K$ & \cite[5.6]{dggw2008}, \cite[Rem.~4.2]{schueth2019}\\
$t^2$ & orbifold, any metric & $K^2$, $\Delta_gK$ (Theorem 4.1 of \cite{schueth2019}) & \cite[Thm~4.1]{schueth2019}\\
all $t^l$ & orbifold, constant curvature $\kappa$ & $\kappa^l$: $\kappa^lp_l(m)/m$ & \cite{ucar2017}, Paper A (bl)\\
$t^{1/2}$, $t\log t$ & non-orbifold cone, $f''(0)\ne0$ & $f''(0)$, irrational in $f'(0)$ & \cite[(3), Thm~4.1]{schueth2025}\\
$t^0$ & any cone $dr^2+f^2d\theta^2$ & $\frac1{12}(\frac1{f'(0)}-f'(0))$ & \cite[(2)]{schueth2025}
\end{tabular}
\end{center}
Schueth's $t^2$ formula is
$a_2(p)=[\frac{m^5-\frac1m}{2520}+\frac{m^3-\frac1m}{720}+\frac{m-\frac1m}{180}]K^2
-[\frac{m^5-\frac1m}{15120}+\frac{m^3-\frac1m}{1440}+\frac{m-\frac1m}{180}]\Delta_gK$; in it the
\emph{derivative term already carries the top power $m^5$}. For rotationally symmetric cones with
$f''(0)=0$ Schueth notes, without proof, that $b_1,b_2$ extend rationally to non-integer $m=1/f'(0)$
\cite[p.~2]{schueth2025}; for $f''(0)\ne0$ half powers and logarithms appear and $b_{1/2}$ is an
irrational function of the cone angle \cite[Abstract, Thm~4.1]{schueth2025}. Orbifold cone points always
have $f''(0)=0$ \cite[p.~2]{schueth2025}.

\subsection*{2. Structure}

\begin{lemma}[Pole order]\label{vc:pole}
Let $\Phi$ be an isometry of a neighbourhood of $p$ with $\Phi(p)=p$, $d\Phi_p$ the rotation by
$\phi\in(0,2\pi)$, $C=C_\phi$. Then $b_l(\Phi)=\sum_{j=0}^{l}C^{-2-2j}T_{l,j}(\phi)$, where the $T_{l,j}$ are
trigonometric polynomials in $\phi$ whose coefficients are universal polynomials in the jet of the
metric at $p$, and
\begin{equation}\label{vc:top}
T_{l,l}(0)=\frac{4^l\,l!}{2\pi}\int_{|u|=1}\bigl[v^{2l}\bigr]\,\rho_u'(v)\,du ,
\end{equation}
$\rho_u$ the inverse function of $\ell_u(r)=|J_u(r)|$, $J_u$ the Jacobi field along $r\mapsto\exp_p(ru)$ with
$J_u(0)=0$, $J_u'(0)$ a unit vector orthogonal to $u$.
\end{lemma}

\begin{proof}
We follow Schueth's proof of \cite[Thm~3.7]{schueth2019}, keeping all orders. With
$d_u(r)=\operatorname{dist}(\exp_p ru,\exp_p rD_\phi u)$ and the substitution $z=d_u(r)/\sqrt t$, $b_l(\Phi)$ is the
$t^l$ coefficient of
\[
\int_{S^1}\int_0^\infty\frac{e^{-z^2/4}}{4\pi t}\Bigl(\sum_iu_i(\exp_p ru,\exp_p rD_\phi u)\,t^i\Bigr)
\sqrt t\,\ell_u(r)\,(d_u^{-1})'(z\sqrt t)\,dz\,du,\qquad r=d_u^{-1}(z\sqrt t)
\]
\cite[(15),(16)]{schueth2019}, the integrand expanded in powers of $\sqrt t$.
(a) The Taylor coefficients of $\operatorname{dist}^2(\exp_px,\exp_py)$ at $(0,0)$ are universal polynomials
in $(x,y)$ and the jet, so $d_u(r)^2=\sum_{n\ge2}\delta_{u,n}(\phi)r^n$ with trigonometric polynomials
$\delta_{u,n}$ and $\delta_{u,2}=C^2$. Each $\delta_{u,n}$ vanishes to second order at $\phi=0$, since
$\operatorname{dist}(\exp_pru,\exp_prD_\phi u)=O(\phi)$; a trigonometric polynomial
$A(\cos\phi)+\sin\phi\,B(\cos\phi)$ with a double zero at $0$ has $A(1)=B(1)=0$, so it is divisible by
$1-\cos\phi=C^2/2$. Hence $d_u(r)=Cr\,W_u(r)^{1/2}$, $W_u=1+\sum_{n\ge1}w_{u,n}(\phi)r^n$, $w_{u,n}$
trigonometric polynomials. (b) Inverting $s=d_u(r)$ formally gives $r=(s/C)V_u(s/C)$ with
$V_u=1+\sum_nv_{u,n}(\phi)(s/C)^n$, so
$\sqrt t\,\ell_u(r)(d_u^{-1})'(s)=\frac{zt}{C^2}\sum_{n\ge0}\gamma_{u,n}(\phi)\bigl(\frac{z\sqrt t}{C}\bigr)^n$,
$\gamma_{u,0}=1$, all $\gamma_{u,n}$ trigonometric polynomials. (c) Likewise
$u_i(\exp_pru,\exp_prD_\phi u)=\sum_n\mu_{i,u,n}(\phi)(z\sqrt t/C)^n$, and since
$u_0(q,q')-1=O(\operatorname{dist}(q,q')^2)$ \cite[(6)]{schueth2019}, every $\mu_{0,u,n}$ with $n\ge1$ is divisible by
$C^2$. (d) A term of order $t^l$ is $t^i(z\sqrt t/C)^N$ with $i+N/2=l$ times $zt/C^2\cdot(4\pi t)^{-1}$ and a
trigonometric polynomial; its pole order in $C$ is $\le2+N\le2+2l$, with equality only if $i=0$ and the
factor from $u_0$ is $1$. Integrating $\int_0^\infty e^{-z^2/4}z^{N+1}dz=2^{N+1}\Gamma(\frac N2+1)$ gives the
first claim (half-integer powers of $t$ cancel, the expansion of $b_l$ having integer powers
\cite[Prop.~4.11]{dggw2008}). (e) The $C^{-2-2l}$ coefficient is
$\frac{1}{4\pi}2^{2l+1}l!\int_{S^1}\gamma_{u,2l}(\phi)du$. As $\phi\to0$, $\delta_{u,n}(\phi)/C^2$ tends to the
$r^n$ coefficient of $\ell_u(r)^2$, because $\partial_\phi\exp_p(rD_\phi u)|_{\phi=0}$ is the Jacobi field of
length $\ell_u(r)$ and $C/\phi\to1$; so $W_u(r)r^2\to\ell_u(r)^2$, $V_u(v)v\to\rho_u(v)$ and
$\gamma_{u,2l}(0)=[v^{2l}]\rho_u'(v)$, which is \eqref{vc:top}.
\end{proof}

\begin{lemma}[Radial reduction]\label{vc:radial}
Let $P$ be a rotation-invariant polynomial of weight $2l$ in the jet of $K$ at $p$ (equivalently in the
Taylor coefficients of $K$ in normal coordinates). If $p$ is a cone point of order $m\ge l-1$, then $P$
takes at $p$ the value it takes on the radial part of the jet (its rotation average), which is the jet of a
rotationally symmetric metric $dr^2+f(r)^2d\theta^2$.
\end{lemma}

\begin{proof}
The jet is $\mathbb Z_m$-invariant, so its non-radial Fourier modes have frequencies $n\equiv0\pmod m$,
$n\ne0$, and occur in Taylor degree $\ge|n|\ge m$, i.e.\ at weight $\ge m+2$. Write $J=J_0+J_{\ne}$. The
part of $P(J)$ linear in $J_{\ne}$ is a rotation-invariant linear form on non-zero frequencies, hence $0$;
the part of degree $\ge2$ in $J_\ne$ has weight $\ge2m+4>2l$, hence is absent. A radial jet of $K$ is the
jet of the rotationally symmetric metric with $f''=-Kf$, $f(0)=0$, $f'(0)=1$.
\end{proof}

\begin{theorem}[Structure of a cone contribution]\label{vc:VC1}
For each $l\ge0$ there are universal rotation-invariant polynomials $\beta_{l,1},\dots,\beta_{l,l+1}$ of
weight $2l$ in $K$ and its covariant derivatives (of order $\le2l-2$) such that:
\begin{enumerate}[label=(\roman*)]
\item for a rotationally symmetric germ, $b_l(\phi)=\sum_{i=1}^{l+1}\beta_{l,i}\,C_\phi^{-2i}$ for every $\phi\in(0,2\pi)$;
\item for every cone point $p$ of order $m\ge\max(2,l-1)$ of every Riemannian orbisurface,
$a_l(p)=\sum_{i=1}^{l+1}\beta_{l,i}(p)\,\Pi_i(m)$ (values on the radial part of the jet);
\item $m\,\Pi_i(m)$ is an even polynomial of degree $2i$ with $m\Pi_i(1)=0$ and leading coefficient
$|B_{2i}|/(2i)!$; $m\Pi_1=\frac{m^2-1}{12}$, $m\Pi_2=\frac{(m^2-1)(m^2+11)}{720}$,
$m\Pi_3=\frac{(m^2-1)(2m^4+23m^2+191)}{60480}$, $m\Pi_4=\frac{(m^2-1)(m^2+11)(3m^4+10m^2+227)}{3628800}$.
\end{enumerate}
So, as at constant curvature, the order-$t^l$ term is a combination of $m^{2l+1},m^{2l-1},\dots,m,1/m$,
one new odd power per order; what changes is the coefficient of each power.
\end{theorem}

\begin{proof}
(i) Every rotation is an isometry of the germ, and so is a reflection, which conjugates $\Phi_\phi$ to
$\Phi_{-\phi}$; so $b_l$ is even in $\phi$. Donnelly's form \cite[4.2]{dggw2008} gives
$b_l=|\det B|\,\tilde b_l(B,\text{jet})$ with $B=(I-D_\phi)^{-1}=\frac12(I+\cot\frac\phi2\,J)$, $J$ the rotation by
$\pi/2$ (the matrix in \cite[5.6]{dggw2008}), $|\det B|=C^{-2}$; hence $b_l=C^{-2}P(\cot\frac\phi2)$ with $P$
a polynomial. By Lemma~\ref{vc:pole}, $C^{2l+2}b_l$ is bounded as $\phi\to0$, so $\deg P\le2l$; $P$ is even,
and $\cot^2\frac\phi2=4C^{-2}-1$. (ii) By Lemma~\ref{vc:radial} applied to the coefficients of $\tilde b_l$
(rotation invariant for fixed $B$, since $B$ commutes with rotations), $b_l(\Phi^j)$ equals (i) at
$\phi=2\pi j/m$ for the radial part of the jet; average over $j$. (iii) Differentiating
$\sum_{j=0}^{m-1}\cot(x+\pi j/m)=m\cot(mx)$ (proof of Paper~A, Lemma~\ref{lem:Phi}) gives
$\sum_{j=1}^{m-1}\csc^2(x+\pi j/m)=m^2\csc^2(mx)-\csc^2x$, whose Taylor coefficients at $x=0$ are
polynomials in $m^2$ vanishing at $m=1$: with $\csc^2y=\sum_nc_ny^{2n-2}$ the $x^{2k}$ coefficient is
$c_{k+1}(m^{2k+2}-1)$. The odd derivatives cancel between $j$ and $m-j$, and
$(d/dx)^{2k}\csc^2x$ is a polynomial in $\csc^2x$ of degree $k+1$ with leading coefficient $(2k+1)!$, so
$\sum_j\csc^{2i}(\pi j/m)$ is obtained triangularly, a polynomial in $m^2$ of degree $i$ vanishing at $m=1$. Its leading coefficient is $2\zeta(2i)/\pi^{2i}=4^i|B_{2i}|/(2i)!$ (the terms $j$ and $m-j$
near the ends dominate). The explicit forms are computed exactly in \texttt{twisted\_mp.py} (reciprocal
power sums of the roots of $T_m(1-S/2)-1$, Newton's identities) and checked numerically.
\end{proof}

\begin{theorem}[The coefficient of the new power]\label{vc:VC2}
$\beta_{l,l+1}$ is given by \eqref{vc:top}; for a rotationally symmetric germ
$\beta_{l,l+1}=4^l\,l!\,[v^{2l}](f^{-1})'(v)$. Hence:
\begin{enumerate}[label=(\roman*)]
\item $\beta_{l,l+1}=\dfrac{2}{(l-1)!}(-\Delta_g)^{l-1}K(p)+N_l$, $N_l$ a polynomial each of whose monomials
has degree $\ge2$ in $K$ and its derivatives ($l\ge1$);
\item at constant curvature $\kappa$, $\beta_{l,l+1}=\frac{(2l)!}{l!}\kappa^l$;
\item $\beta_{1,2}=2K$, $\beta_{2,3}=12K^2-2\Delta_gK$, $\beta_{3,4}=120K^3-42K\Delta_gK+\Delta_g^2K$,
$\beta_{4,5}=1680K^4-904K^2\Delta_gK+\frac{98}3K\Delta_g^2K+\frac{124}3(\Delta_gK)^2-\frac13\Delta_g^3K$
(radial jets; for $l=4$ this is the value at cone points of order $m\ge3$);
\item for $m\ge\max(2,l-1)$ the coefficient of $m^{2l+1}$ in $a_l(p)$ is
$\dfrac{|B_{2l+2}|}{(2l+2)!}\,\beta_{l,l+1}(p)$.
\end{enumerate}
\end{theorem}

\begin{proof}
The first sentence is \eqref{vc:top}, with $\ell_u=f$ in the symmetric case. (ii) $f=\sin(\sqrt\kappa
r)/\sqrt\kappa$, $(f^{-1})'(v)=(1-\kappa v^2)^{-1/2}=\sum_a\binom{2a}a(\kappa v^2/4)^a$. (i) To first order in
$K=\sum_jk_jr^{2j}$, $f=r-\int_0^r(r-s)K(s)s\,ds$, so $f\ni-k_{l-1}r^{2l+1}/(2l(2l+1))$,
$\rho\ni k_{l-1}v^{2l+1}/(2l(2l+1))$, $\rho'\ni k_{l-1}v^{2l}/(2l)$, and
$(-\Delta_g)^{l-1}K(p)=4^{l-1}((l-1)!)^2k_{l-1}+(\text{degree}\ge2)$ since
$(\partial_1^2+\partial_2^2)^{n}r^{2n}=4^n(n!)^2$. In general (non-radial) the linear part of the
$u$-average of the order-$(2l-2)$ Taylor coefficient of $K$ is the same trace. (iii) Exact computation,
\texttt{top\_coefficient.py} ($l\le7$), with $l=1,2$ equal to Donnelly's $b_1=2K/C^4$
\cite[Rem.~3.2]{schueth2019} and the $C^{-6}$ part of \cite[Thm~3.7]{schueth2019}, and $l=3$ equal to the
$C^{-8}$ part of $b_3$ computed independently in Theorem~\ref{vc:VC3}. (iv) Theorem~\ref{vc:VC1}(iii) with
$i=l+1$.
\end{proof}

\begin{corollary}[What fails]\label{vc:fail}
For every $l\ge2$ the top power $m^{2l+1}$ of the order-$t^l$ cone term has coefficient
$\frac{|B_{2l+2}|}{(2l+2)!}\bigl(\frac{(2l)!}{l!}K^l+\cdots+\frac{2}{(l-1)!}(-\Delta_g)^{l-1}K\bigr)$, which involves
$\Delta_g^{l-1}K(p)$ with non-zero coefficient. The ansatz ``leading $m$-dependence $K(p)^lp_l(m)/m$, derivative
terms with lower powers of $m$'' therefore holds for $l=0,1$ only. At $l=2$ this is visible in
\cite[Thm~4.1]{schueth2019}: $-\frac1{15120}m^5\Delta_gK$.
\end{corollary}

The failure is not confined to the top coefficient. To first order in the curvature, the derivative of
$K$ enters the order-$t^l$ cone term \emph{only} through the new power sum, for every order $m\ge2$:

\begin{theorem}[The linear part of a cone contribution]\label{vc:VC7}
Let $l\ge1$.
\begin{enumerate}[label=(\roman*)]
\item Let $\Phi$ be an isometry of a neighbourhood of $p$ with $\Phi(p)=p$, $d\Phi_p$ the rotation by
$\phi\in(0,2\pi)$, $C=C_\phi$. Then
$b_l(\Phi)=\dfrac{2}{(l-1)!}\,C^{-2l-2}\,(-\Delta_g)^{l-1}K(p)+N_l$, where $N_l$ is a polynomial in the jet of $K$
at $p$ (coefficients depending on $\phi$) each of whose monomials has degree $\ge2$.
\item For every cone point $p$ of order $m\ge2$ of every Riemannian orbisurface,
\[
a_l(p)=\frac{2}{(l-1)!}\,\Pi_{l+1}(m)\,(-\Delta_g)^{l-1}K(p)+Q_{l,m}(p),
\]
$Q_{l,m}$ a polynomial in the jet of $K$ at $p$ each of whose monomials has degree $\ge2$.
\end{enumerate}
\end{theorem}

\begin{proof}
(i) By Donnelly's form \cite[4.2]{dggw2008}, $b_l(\Phi)=C^{-2}\tilde b_l$ with $\tilde b_l$ a polynomial in $B$ and
in $K$ and its covariant derivatives at $p$, invariant under the rotations (which commute with $B$). The
monomials of weight $2l\ge2$ that are linear in the jet are contractions of $\nabla^{2l-2}K(p)$ with tensors
built from $B$, $g$ and the area form; as $\nabla^{2l-2}K$ is symmetric up to terms quadratic in the jet, and
the rotation-invariant linear forms on $\operatorname{Sym}^{2l-2}\mathbb R^2$ are the multiples of the full
trace, the linear part of $b_l$ is $\lambda_l(\phi)(-\Delta_g)^{l-1}K(p)$ for a function $\lambda_l$ alone. To
find $\lambda_l$ it suffices to evaluate one family with a non-zero linear part: the
germ $g_\varepsilon=e^{2\varepsilon\psi}|dx|^2$, $\psi=|x|^{2l}$, for which $R_\phi$ is an isometry. Since $g_0$ is
flat and $b_l$ is a polynomial in the jet vanishing on the flat jet, $\frac{d}{d\varepsilon}|_0b_l$ is the linear part
evaluated on $\frac{d}{d\varepsilon}|_0(\text{jet})$. (In general isothermal coordinates make $\Phi$ linear: the Riemann
map of a $\Phi$-invariant geodesic disc onto the unit disc, $p\mapsto0$, conjugates $\Phi$ to a conformal
automorphism fixing $0$, hence to $R_\phi$ by Schwarz's lemma; this is how \texttt{linear\_part.py} tests arbitrary
$\psi$.) As
$\Delta_g=e^{-2\varepsilon\psi}\Delta_0$, $\Delta_0=-(\partial_1^2+\partial_2^2)$, and $R_\phi^*$ commutes with $\Delta_0$, Duhamel's formula
gives, up to $O(t^\infty)$ from a cut-off,
\[
\frac{d}{d\varepsilon}\Big|_0\operatorname{tr}\bigl(R_\phi^*e^{-t\Delta_g}\bigr)
=2t\operatorname{tr}\bigl(R_\phi^*\psi\Delta_0e^{-t\Delta_0}\bigr)
=2t\int\psi(x)\,h_t(Cx)\Bigl(\frac1t-\frac{C^2|x|^2}{4t^2}\Bigr)dx,
\]
$h_t(z)=e^{-|z|^2/4t}/4\pi t$, because $|R_\phi x-x|=C|x|$. The integrand depends on $x$ only through
$\psi$ times a radial function, so only the radial part of $\psi$ contributes; for $\psi=r^{2k}$ the
substitution $s=C^2r^2/4t$ gives $-2k\,k!\,4^kt^kC^{-2k-2}$, a single power of $t$ and of $C$; put $k=l$. On the other hand the linear part of
$(-\Delta_g)^{l-1}K(0)$ is $(-1)^{l-1}\Delta_0^l\psi(0)$ (since $K=e^{-2\varepsilon\psi}\Delta_0(\varepsilon\psi)$), and
$\Delta_0^lr^{2k}(0)=-\delta_{kl}(-1)^{l-1}4^l(l!)^2$. Comparing at $k=l$: $\lambda_l=-2l\,l!\,4^lC^{-2l-2}/(-4^l(l!)^2)
=\frac2{(l-1)!}C^{-2l-2}$. (Any $\varepsilon$-linear jet is realised by some $\psi$, and the linear part does
not depend on the coordinates in which the jet is written, the change of coordinates being the identity
to first order in the curvature.)
(ii) Average (i) over $\Phi^j$, $j=1,\dots,m-1$; no restriction on $m$ is needed, because the linear part
sees no non-radial jet.
\texttt{linear\_part.py} asserts the computation monomial by monomial ($\psi=x^ay^b$, $a+b\le14$,
$l\le7$, radial and non-radial alike), and checks it against Donnelly's $b_1$, the $\Delta_gK$ term of
\cite[Thm~3.7]{schueth2019}, Theorem~\ref{vc:VC3}, and the attacker's $b_4$; after averaging, against the
$K$ term of \cite[Rem.~4.2]{schueth2019} ($2\Pi_2$), the $\Delta_gK$ term of \cite[Thm~4.1]{schueth2019}
($-2\Pi_3$), and $D(m)=\Pi_4(m)$.
\end{proof}

So the ansatz of the question is inverted, not merely perturbed: at linear order the highest derivative
$(-\Delta_g)^{l-1}K(p)$ is the coefficient of the full polynomial $\Pi_{l+1}(m)$, which carries the new top power
$m^{2l+1}$, while $K(p)^l$ and every other non-linear monomial appear in the coefficients of
$\Pi_1,\dots,\Pi_{l+1}$. Summed over cone points, the order-$t^l$ cone term is, to first order,
$\frac{2}{(l-1)!}\sum_i(-\Delta_g)^{l-1}K(p_i)\,\Pi_{l+1}(m_i)$: a power sum \emph{weighted by an independent
real number at each cone point}, which is the precise sense in which the Prouhet--Tarry--Escott structure is
lost.

\begin{theorem}[The $t^3$ cone coefficient]\label{vc:VC3}
For a rotationally symmetric germ,
$b_3(\phi)=\bigl(\frac{120}{C^8}-\frac{32}{C^6}+\frac4{3C^4}\bigr)K^3+\bigl(-\frac{42}{C^8}+\frac8{C^6}\bigr)K\Delta_gK+\frac1{C^8}\Delta_g^2K$.
For every cone point of order $m\ge2$ of every Riemannian orbisurface,
\[
a_3(p)=\tfrac{(m^2-1)(m^2+3)(3m^4+2m^2+19)}{30240\,m}K^3
-\tfrac{(m^2-1)(7m^6+47m^4+173m^2+733)}{201600\,m}K\Delta_gK
+\tfrac{(m^2-1)(m^2+11)(3m^4+10m^2+227)}{3628800\,m}\Delta_g^2K .
\]
\end{theorem}

\begin{proof}
\texttt{twisted\_mp.py}: for $K=k_0+k_1r^2+k_2r^4$ it solves the Hamilton--Jacobi equation for
$\sigma=\operatorname{dist}^2/2$ and the transport equations
$\sigma^a\nabla_aU_i+(i+\frac12(\nabla^2\sigma-2))U_i=\nabla^2U_{i-1}$, $U_0(x,x)=1$, exactly as power series in
normal coordinates (residuals asserted zero), and evaluates
$\int_0^\infty2\pi f(r)\,(4\pi t)^{-1}e^{-\sigma/2t}\sum_iU_it^i\,dr$ at the points $(r,0)$, $R_\phi(r,0)$ by
Laplace's method term by term (this is Schueth's computation in the variable $r$ instead of $z$).
Checks, all asserted: $U_1(p,p)=K/3$, $U_2(p,p)=K^2/15-\Delta_gK/15$, $U_0$ and $U_1$ along a ray agree with
\cite[(6),(8)]{schueth2019}; $b_0=C^{-2}$, $b_1=2KC^{-4}$, $b_2$ equals \cite[Thm~3.7]{schueth2019}; after
averaging, $a_0,a_1,a_2$ equal \cite[Rem.~4.2, Thm~4.1]{schueth2019}; the $K^3$ part of $a_3$ equals
U\c{c}ar's $p_3(m)/m$; the $C^{-8}$ part equals Theorem~\ref{vc:VC2}(iii). For general cone points: the
rotation-invariant weight-$6$ polynomials are spanned by $K^3$, $K\Delta_gK$, $|\nabla K|^2$, $\Delta_g^2K$
(other contractions of $\nabla^4K$ differ from $\Delta_g^2K$ by these), $\nabla K(p)=0$ at a cone point, and
Lemma~\ref{vc:radial} applies for $m\ge2=l-1$; the radial family realises all values of
$(K,\Delta_gK,\Delta_g^2K)=(k_0,-4k_1,64k_2-\frac83k_0k_1)$.
\end{proof}

\subsection*{3. An extension}

\begin{theorem}[Curvature flat to finite order at the cone points]\label{vc:VC4}
Fix $\kappa\ne0$ and $L\ge2$. Let $(\Orb,g)$, $(\Orb',g')$ have signatures $(\gamma;m)$, $(\gamma';m')$,
and suppose that at every cone point of both, $K(p)=\kappa$ and $\nabla^jK(p)=0$ for $1\le j\le2L-6$.
Suppose $S_l(g)=S_l(g')$ for $-1\le l\le L-2$. Then $\chi^{\rm orb}(\Orb)=\chi^{\rm orb}(\Orb')$, and
$a_l(g)=a_l(g')$ for $-1\le l\le L-2$ (that is, $c_1,\dots,c_L$ agree) if and only if
$\Psi_k(m)=\Psi_k(m')$ for $1\le k\le L-1$. Hence the reduction of Paper~A, Lemma~\ref{lem:sigdata}, holds
verbatim in this class, and with it every conclusion Paper~A draws from that lemma, with $\Area$ replaced
by $-2\pi\chi^{\rm orb}$.
The same holds for $L\le5$ if instead all cone points of $\Orb$ and $\Orb'$ have the same
$J=(K,\Delta_gK,\Delta_g^2K)(p)$ and $\beta_{l,l+1}(J)\ne0$ for $1\le l\le L-2$. For $\kappa=0$ and $L\ge3$
the ``only if'' fails.
\end{theorem}

\begin{proof}
For $l\le L-2$, $a_l(p)$ is a polynomial of weight $2l$ in the jet, so it involves $\nabla^jK(p)$ only for
$j\le2l-2\le2L-6$; these coincide with the constant-curvature values, and locality
\cite[4.1]{dggw2008} gives $a_l(p)=\kappa^lp_l(m)/m$ (U\c{c}ar; Paper~A, \eqref{eq:bl} with $\kappa^l$ for
$(-1)^l$). $S_0=\chi^{\rm orb}/6$ gives the first claim. For $0\le l\le L-2$,
$a_l(g)-a_l(g')=\kappa^l(C_l(m)-C_l(m'))$ in Paper~A's notation, and $\kappa\ne0$; Paper~A,
Lemma~\ref{lem:triangular}, turns this into $\Psi_k$, $k\le L-1$; $a_{-1}=S_{-1}$. For a common $J$,
$\sum_ia_l(p_i)=\sum_{d\le l+1}\beta_{l,d}(J)\sum_i\Pi_d(m_i)$ (Theorems~\ref{vc:VC1}, \ref{vc:VC3}; explicit for
$l\le3$) is triangular with diagonal $\beta_{l,l+1}(J)$, and $m\Pi_d(m)$, $d\le n$, span the same space as
$m^{2d}-1$, $d\le n$ (Theorem~\ref{vc:VC1}(iii)). For $\kappa=0$, $a_l(p)=0$ for $1\le l\le L-2$, and
Theorem~\ref{vc:VC6} gives pairs with all $a_l$ and all $S_l$ equal but $\Psi_2(m)\ne\Psi_2(m')$.
\end{proof}

The smooth parts $S_l$ are not spectral data, and Theorems~\ref{vc:VC5}--\ref{vc:VC6} show that this
hypothesis cannot be dropped.

\subsection*{4. Obstructions}

\begin{lemma}[Second variation of the smooth invariants]\label{vc:quad}
Let $\psi\in C_c^\infty(\mathbb R^2)$ and $g_\lambda=e^{2\lambda\psi}|dx|^2$. For $l\ge1$,
\[
G_l(\lambda\psi):=\frac1{4\pi}\int u_{l+1}(g_\lambda)\,dA_\lambda
=\lambda^2\,\frac{F_{l+1}}{2\pi}\int\psi\,\Delta_0^{\,l+1}\psi\,dx+O(\lambda^3),\qquad
F_n=\frac{(-1)^n\,n(n-1)\,n!}{(2n+1)!}\ne0\ (n\ge2).
\]
\end{lemma}

\begin{proof}
Here $\Delta_0=-(\partial_1^2+\partial_2^2)\ge0$. Put $\psi$ in a large flat torus; $\Delta_{g_\lambda}=e^{-2\lambda\psi}\Delta_0=\Delta_0+w\Delta_0$,
$w=e^{-2\lambda\psi}-1$. The heat expansion and its remainder depend smoothly on $\lambda$, so the
$\lambda^2$ coefficient of $G_l$ is the $t^l$ coefficient of the $\lambda^2$ part of
$\operatorname{tr}e^{-t\Delta_{g_\lambda}}$. By Duhamel, that part is
$-t\operatorname{tr}(w_2\Delta_0e^{-t\Delta_0})+\frac{t^2}2\int_0^1\operatorname{tr}(w_1\Delta_0e^{-(1-\alpha)t\Delta_0}w_1\Delta_0e^{-\alpha t\Delta_0})d\alpha$,
$w_1=-2\lambda\psi$, $w_2=2\lambda^2\psi^2$. The first term is $-(4\pi t)^{-1}\int w_2$ up to $O(t^\infty)$. In
Fourier variables (lattice sums equal the integrals up to $O(t^\infty)$), completing the square
$(1-\alpha)|\eta|^2+\alpha|\eta+\zeta|^2=|\eta+\alpha\zeta|^2+\alpha(1-\alpha)|\zeta|^2$ and using the Gaussian
moment $\mathbb E|\mu+A|^2|\mu+B|^2=2t^{-2}+(|A|^2+|B|^2+2A\cdot B)t^{-1}+|A|^2|B|^2$, the second term is
$\frac{1}{2\pi t}\cdot\frac1{(2\pi)^2}\int|\hat\psi(\zeta)|^2F(t|\zeta|^2)\,d\zeta\cdot\lambda^2$ with
$F(x)=\int_0^1[2+(1-2\alpha)^2x+\alpha^2(1-\alpha)^2x^2]e^{-\alpha(1-\alpha)x}d\alpha=\sum_nF_nx^n$; the
integral $\int_0^1(\alpha(1-\alpha))^nd\alpha=(n!)^2/(2n+1)!$ gives $F_n$. Plancherel gives the claim.
Checks (\texttt{quadratic\_form.py}): $F_1=0$ (no $t^0$ term: $\int K\,dA=0$ for a bump); at $l=1$ the
formula is $\frac1{60\pi}\int(\Delta_0\psi)^2\lambda^2=\frac1{4\pi}\int\frac{K^2}{15}$, Berger's $u_2$
\cite[(7)]{schueth2019}; at $t^{-1}$ the two terms add to the $\lambda^2$ part of $\Area/4\pi$.
\end{proof}

\begin{lemma}[Generalised Vandermonde]\label{vc:vdm}
For distinct reals $e_1,\dots,e_N$ and distinct $\varepsilon_1,\dots,\varepsilon_N>0$,
$\det(\varepsilon_k^{e_i})\ne0$.
\end{lemma}

\begin{proof}
A non-zero $\sum_ic_i\varepsilon^{e_i}$ has at most $N-1$ positive zeros: divide by the smallest power and
differentiate (Rolle, induction on $N$).
\end{proof}

\begin{theorem}[Finite order: only $c_2$ survives]\label{vc:VC5}
Let $\sigma=(\gamma;m_1,\dots,m_n)$ and $\sigma'=(\gamma';m'_1,\dots,m'_{n'})$ be signatures of closed
orientable cone-point orbifolds, $\kappa\in\mathbb R$ and $L\ge0$. There are orbifold metrics $g$ on
$\Orb_\sigma$ and $g'$ on $\Orb_{\sigma'}$, each of constant curvature $\kappa$ near every cone point, with
$a_l(g)=a_l(g')$ for all $-1\le l\le L$, if and only if
\begin{equation}\label{vc:a0}
\frac{2-2\gamma}6+\sum_i\frac{(m_i-1)^2}{12m_i}=\frac{2-2\gamma'}6+\sum_j\frac{(m'_j-1)^2}{12m'_j}.
\end{equation}
For example $S^2(2,2,2,2,2,2,2,2)$ and $S^2(3,3,3)$ (both sides $2/3$).
\end{theorem}

\begin{proof}
Only if: $a_0=\chi^{\rm orb}/6+\sum_i\frac{m_i^2-1}{12m_i}$ for every metric \cite[(5.7)]{dggw2008}, and this is
the left side of \eqref{vc:a0}.
If: (1) Choose $g_0$ on $\Orb_\sigma$ of constant curvature $\kappa$ on disjoint discs $U_i$ about the cone
points (quotients of small geodesic discs of the space form by rotations of order $m_i$), flat on two
further disjoint discs $\Delta_1,\Delta_2$ (convex combination of metrics). Inside $\Delta_2$ replace the flat
metric by a rotationally symmetric ``finger'' containing a flat cylinder $S^1_\rho\times[0,T]$, all of the
finger except the cylinder independent of $T$: the area grows by $2\pi\rho T$ and no $S_l$ changes.
Do the same on $\Orb_{\sigma'}$. By locality, every metric that agrees with $g_0$ on $\bigcup U_i$ has
cone parts $C_l=\sum_i\kappa^lp_l(m_i)/m_i$. (2) Fix $\psi\in C_c^\infty(B_1)$, $\psi\ne0$, and by
Lemma~\ref{vc:quad} $\lambda\ne0$ with $G_l:=G_l(\lambda\psi)\ne0$ for $1\le l\le L$. The
$\varepsilon$-gadget is the metric $e^{2\lambda\psi((x-x_0)/\varepsilon)}|dx|^2$ in a disc of radius $\varepsilon$
inside $\Delta_1$; it changes $S_l$ by $\varepsilon^{-2l}G_l$ (weights: $S_l(c^2h)=c^{-2l}S_l(h)$) and nothing
else but the area. Gadgets in disjoint discs add. (3) The map
$V(\varepsilon_1,\dots,\varepsilon_L)=\sum_k(\varepsilon_k^{-2}G_1,\dots,\varepsilon_k^{-2L}G_L)$ has Jacobian
$\prod_l(-2lG_l)\cdot\det(\varepsilon_k^{-2l-1})\ne0$ at distinct scales (Lemma~\ref{vc:vdm}), so its image
contains a ball $B(w_0,\rho_0)$. (4) Let $E_l=S_l(g'_0)+C'_l-S_l(g_0)-C_l$ and
$\Sigma_s=\operatorname{diag}(s^{-2},\dots,s^{-2L})$. For $s\in(0,1]$ small,
$\Sigma_s^{-1}E\in B(0,\rho_0)$; take gadget scales $X$ on $\Orb_\sigma$ with $V(X)=w_0+\Sigma_s^{-1}E$ and $Y$ on
$\Orb_{\sigma'}$ with $V(Y)=w_0$, and shrink all gadgets by $s$ about their centres: the changes become
$\Sigma_sV(X)$ and $\Sigma_sV(Y)$, whose difference is $E$. Now $S_l(g)+C_l=S_l(g')+C'_l$ for
$1\le l\le L$. (5) Lengthen one finger to equalise the areas ($a_{-1}$); $a_0$ agree by \eqref{vc:a0}.
\end{proof}

\begin{theorem}[All orders: explicit metrics, flat at the cone points]\label{vc:VC6}
Let $m=(m_1,\dots,m_n)$, $m'=(m'_1,\dots,m'_{n'})$ be multisets of integers $\ge2$ with
\begin{equation}\label{vc:vc6}
\sum_i\Bigl(1-\frac1{m_i}\Bigr)=\sum_j\Bigl(1-\frac1{m'_j}\Bigr)>1,\qquad
\sum_i\Bigl(m_i-\frac1{m_i}\Bigr)=\sum_j\Bigl(m'_j-\frac1{m'_j}\Bigr).
\end{equation}
For every $\gamma\ge0$ there are orbifold metrics on $\Orb_{(\gamma;m)}$ and $\Orb_{(\gamma;m')}$, flat near every
cone point, whose heat traces have the same asymptotic expansion to all orders. This applies to
$\Orb(2,8,8)$ and $\Orb(3,3,12)$ (Paper~A, Theorem~\ref{thm:intro-rigid}(ii)), to $(2k,8k,8k)$ and
$(3k,3k,12k)$, and to $(5,5,5)$ and $(2,2,2,10)$.
\end{theorem}

\begin{proof}
Let $\alpha_i=1/m_i$ and $\alpha_0=\sum_i(1-\alpha_i)-1>0$, so $\sum_{i=0}^n(\alpha_i-1)=-2$. On $\hat{\mathbb C}$ with
distinct $z_0,\dots,z_n\in\mathbb C$, the metric $|\omega|$, $\omega=\prod_{i=0}^n(z-z_i)^{\alpha_i-1}dz$, is flat; at
$\infty$, in $w=1/z$, $|\omega|=|w|^{-\sum(\alpha_i-1)-2}\prod_i|1-z_iw|^{\alpha_i-1}|dw|$ is smooth. Near $z_i$ there
is a holomorphic coordinate $\zeta$ with $\omega=\zeta^{\alpha_i-1}d\zeta$; for $\alpha_i=1/m_i$ the chart
$\zeta=\xi^{m_i}$ turns $|\zeta|^{1/m_i-1}|d\zeta|$ into $m_i|d\xi|$, a flat orbifold cone point of order $m_i$; near
$z_0$, with $\varrho=|\zeta|^{\alpha_0}/\alpha_0$, the metric is the exact flat cone $d\varrho^2+\alpha_0^2\varrho^2d\theta^2$.
Do the same for $m'$ with the same $\alpha_0$, scale both so that the cone neighbourhoods have radius
$>\varrho_1$, remove $\{\varrho<\varrho_1\}$ from both, and glue in the same smooth surface $P_T$ with one
boundary circle: genus $\gamma$, a collar isometric to $\{\varrho_1-\delta<\varrho\le\varrho_1\}$ of that cone, and a flat
cylinder of length $T$ (e.g.\ a surface of revolution $d\varrho^2+F(\varrho)^2d\theta^2$, $F'(0)=1$,
$F=\alpha_0\varrho$ near the collar, with $\gamma$ handles attached in a flat part). The result is a smooth
Riemannian orbifold of signature $(\gamma;m)$, resp.\ $(\gamma;m')$. For $l\ge1$: $u_{l+1}=0$ on flat regions, so
$S_l=\frac1{4\pi}\int_{P_T}u_{l+1}$, the same for both and independent of $T$; and a flat cone point
contributes $a_l(p)=0$ (each monomial of $b_l$ contains a curvature factor). For $l=0$, $\chi^{\rm orb}$ and
$\sum_i\frac{m_i^2-1}{12m_i}=\frac1{12}\sum_i(m_i-\frac1{m_i})$ agree by \eqref{vc:vc6}. For $l=-1$, choose $T,T'$.
\end{proof}

\begin{remark}[What is open]
(a) The lower coefficients $\beta_{l,i}$, $i\le l$, for $l\ge4$ (Theorem~\ref{vc:VC1} gives their existence and
Theorem~\ref{vc:VC2} the top one), and, for $m\le l-2$, the terms quadratic in non-radial jets; these do
occur: at $l=4$, $m=2$ the adversarial check (\texttt{attack-log.md}) finds
$b_4(\pi)=\frac7{480}|\nabla^2K-\frac12\operatorname{tr}\nabla^2K\,g|^2+(\text{radial part})$.
(b) Whether \eqref{vc:a0} alone suffices for agreement to all orders (Theorem~\ref{vc:VC6} needs both
conditions of \eqref{vc:vc6}, i.e.\ equal $\chi^{\rm orb}$ as well). (c) Whether any of these pairs can be
made isospectral. (d) Schueth's conjecture for corners of arbitrary angle \cite[Conj.~5.5]{schueth2019} is
not addressed.
\end{remark}
